% Master CV example: Isaac Newton
\documentclass[11pt,a4paper]{article}
\usepackage[margin=1.7cm]{geometry}
\usepackage{hyperref}
\usepackage{enumitem}
\setlist[itemize]{leftmargin=*,nosep}
\begin{document}
\begin{center}
{\LARGE \textbf{Isaac Newton}}\\
Mathematical Scientist | Experimental Researcher | Technical Innovator\\
Woolsthorpe, Lincolnshire, England\\
\end{center}
\section*{Professional Profile}
Mathematical scientist and experimental researcher with extensive experience developing theoretical models, investigating physical phenomena, designing instruments, and translating complex observations into practical methods. Strong analytical, problem-solving, and technical communication skills, with experience in scientific leadership, operational management, and the improvement of technical processes.
\section*{Professional Experience}
\textbf{Master of the Royal Mint} \hfill 1699--1727\\
Royal Mint, London
\begin{itemize}
\item Oversaw technical and operational activities within a major national manufacturing institution.
\item Investigated counterfeit currency and developed evidence-led methods for identifying and prosecuting offenders.
\item Applied detailed analysis, measurement, and procedural oversight to work requiring accuracy and reliability.
\item Supported the integrity and quality of currency production and associated operations.
\end{itemize}
\textbf{Warden of the Royal Mint} \hfill 1696--1699\\
Royal Mint, London
\begin{itemize}
\item Contributed to the administration and technical oversight of currency production.
\item Conducted investigations into coinage, manufacturing practices, and counterfeiting.
\end{itemize}
\textbf{Lucasian Professor of Mathematics} \hfill 1669--1701\\
University of Cambridge
\begin{itemize}
\item Conducted original research in mathematics, mechanics, optics, and the physical sciences.
\item Developed mathematical methods for describing motion and physical systems.
\item Designed and built a reflecting telescope, applying optical theory to instrument development.
\item Communicated complex technical ideas through lectures, publications, and scientific correspondence.
\end{itemize}
\section*{Selected Technical Achievements}
\begin{itemize}
\item Formulated laws of motion and universal gravitation, establishing mathematical frameworks for analysing mechanical systems.
\item Developed foundational methods in calculus to address problems involving change, motion, and quantities.
\item Investigated light and colour through controlled experiments with prisms and optical apparatus.
\item Designed and constructed a reflecting telescope to address limitations in refracting telescope design.
\item Published \textit{Philosophi\ae Naturalis Principia Mathematica}, integrating mathematical analysis with physical observations.
\end{itemize}
\section*{Education}
\textbf{Master of Arts} -- University of Cambridge, 1668\\
\textbf{Bachelor of Arts} -- University of Cambridge, 1665
\section*{Technical Skills}
Mathematical modelling; mechanics; optics; experimental design; instrument development; quantitative analysis; technical investigation; process oversight; problem-solving; scientific writing; research leadership.
\section*{Leadership and Professional Service}
\begin{itemize}
\item President of the Royal Society, 1703--1727.
\item Knighted in 1705 for services to the Crown.
\end{itemize}
\end{document}
